\chapter{总结与展望}

\section{项目总结}

数智游踪已围绕A5“景区导览服务AI数字人”赛题形成\textbf{完整产品闭环}：游客端以数字人为统一入口，支持文本与语音问答、景点讲解、个性路线与满意度反馈；管理端覆盖身份权限、知识维护、数字人配置、会话质量与运营大屏。系统完成简体中文、英文、韩文、繁体中文和日文的\textbf{全局界面国际化}，并以讯飞云端流式数字人与Live2D本地模型形成\textbf{配置驱动的双引擎体系}；底层以本地知识和引用约束事实边界，并通过SQLite、OpenAPI、自动化测试与单机同源部署保证\textbf{工程可复现性}。

项目价值不依赖单一模型能力，而来自“资料入库—熟悉语言操作—有据回答—数字人表达—数据沉淀—知识修正”的\textbf{持续循环}。50条标准问答、20条检索用例、数字人配置测试与五语言生产构建共同证明，当前实现既可支撑比赛演示，也具备进入\textbf{单景区试运行}的基础。

\section{后续展望}

后续建设坚持\textbf{先验证价值、再扩大范围}：近期接入GPS或蓝牙信标实现到点触发与弱网缓存，并补充\textbf{经人工审校的多语种景点资料}；中期对接票务、客流与活动排期系统，使路线推荐兼顾\textbf{实时拥挤度与运营目标}；规模化阶段进一步将景点实体、语言包、数字人形象与指标看板抽象为\textbf{景区级配置}，形成\textbf{可复制的多租户能力}，并持续校准模型与交互策略。

\section{最后的话}

感谢指导老师\textbf{陈子淮}在选题方向、系统设计与文档打磨过程中的悉心指导；感谢第十五届“中国软件杯”大赛\textbf{组委会}搭建的竞赛与交流平台，让团队得以在真实赛题中锤炼产品与技术；也感谢\textbf{数智游踪队}三位队友的通力协作，从知识库、算法到前端与数字人各司其职、彼此补位，共同完成了这份作品。
